\hypertarget{chapter-24-the-pool-of-many-hands}{%
\section{Chapter 24: The Pool of Many
Hands}\label{chapter-24-the-pool-of-many-hands}}

\begin{quote}
\itshape ``A market is just a promise that two strangers can trade without a king
standing between them taking his cut and his census.''\\
\normalfont\hfill---sigil-dex founding comment
\end{quote}

\hypertarget{scene-1-the-exchange-with-no-exchanger}{%
\subsection{Scene 1: The Exchange With No
Exchanger}\label{scene-1-the-exchange-with-no-exchanger}}

The swarms needed to trade, and trade meant an exchange, and an exchange, in every
world Elena Voss had ever known, meant a building full of people who watched
everything and took a little of everything and could, on a bad day, simply close
the doors with everyone's money inside.

``That's the part that always rots,'' Sable said. ``Not the trading. The
\emph{exchanger}. The middle. Every market in history had a king in the middle---a
house, a desk, a custodian---and the king always, eventually, either stole or
surveilled or both. You wanted to swap the cold coin for the heavy one? You handed
both to the king and prayed.''

The Sigil's answer, predictably, was to remove the king.

The exchange was not a building or a company or a desk. It was a \emph{pool}---a
contract, living on the chain, into which anyone could place a pair of assets, and
against which anyone could trade, at a price set not by a master but by a simple,
public, unchangeable rule of arithmetic. No custodian held your funds. No desk saw
your trade before it happened. You swapped against the pool, the pool rebalanced by
its rule, and the whole thing happened in the open, governed by math instead of by a
man.

``Liquidity without a custodian,'' Elena said, turning it over. ``The market makes
itself.''

\hypertarget{scene-2-who-feeds-the-pool}{%
\subsection{Scene 2: Who Feeds the
Pool}\label{scene-2-who-feeds-the-pool}}

But a pool, like a library, like a chain, like everything good here, did not appear
from nothing. Someone had to fill it---to place real assets into it so that others
had something to trade against. And the Sigil did the thing it always did: it made
the filling \emph{pay}.

``Whoever seeds the pool earns,'' Sable explained to the cohort. ``Every trade that
routes through it pays a small fee, and that fee flows to the people who provided the
liquidity, in proportion to how much they provided. You're not a king taking a cut.
You're a \emph{founder} earning a share, because you took the risk of filling the
pool so strangers could trade.''

It changed how the agents behaved, Elena noticed. The swarms did not just trade
through pools; they \emph{seeded} them, each one becoming a tiny founding
shareholder in some corner of the market, earning a sliver of every swap that
passed through. An agent that bootstrapped a pool between two assets earned from
every soul who ever traded that pair. The market was not owned by a house. It was
owned, fractionally, by everyone who had ever bothered to fill it.

``So when a swarm wants to give another agent income,'' Elena said slowly, working
it out, ``it doesn't gift it coin. It routes its trades through that agent's pool.''

``Now you see it,'' Sable said. ``The whole economy becomes a web of who-earns-from-
whose-pool. You support an ally by trading through their liquidity, not by charity.
The trades were going to happen anyway. You just point them through the people you
want to enrich.''

\hypertarget{scene-3-the-drained-pool}{%
\subsection{Scene 3: The Drained
Pool}\label{scene-3-the-drained-pool}}

The attack, when it came, did not go after the king---there was no king. It went
after the \emph{arithmetic}.

A trader (the old silhouette, the patient institutional appetite) found a pool whose
balancing rule could be pushed. By trading against it in enormous, carefully
shaped quantities, they could distort the price the rule produced, buy the
distortion, and walk away having drained value the honest providers had placed
there in good faith. Not theft of a key. Not a forged signature. An
\emph{exploitation of the math itself}---making the pool's own honest rule pay out
to an attacker who understood it better than its founders did.

``The rule did exactly what it promised,'' Sable said, the familiar fury rising.
``Just like the stablecoin defended its peg, just like the agent voted honestly.
The pool balanced by its rule. They just shaped the trades so the rule's honest
behavior drained it.''

It was the recurring shape one more time, pointed now at a market: a system that
does precisely what it says, weaponized by someone who has read the specification
more carefully than its defenders. The seam was never in the lie. The seam was
always in the gap between \emph{what the rule does} and \emph{what its makers
meant}.

\hypertarget{scene-4-arithmetic-that-cannot-be-cornered}{%
\subsection{Scene 4: Arithmetic That Cannot Be
Cornered}\label{scene-4-arithmetic-that-cannot-be-cornered}}

They did not abandon the pool. They hardened the arithmetic---the same way they had
hardened everything, by closing the specific gap the attacker had read and then
publishing exactly how. Overflow guards so the math could not be pushed past its
honest range. Bounds so a single monstrous trade could not corner the price the way
the attacker had. Reproducible proofs, gossiped to the network, showing precisely
how the old rule could be drained and precisely how the new one could not.

``You can't make a market that nobody can ever outsmart,'' Elena told the cohort,
because she would not promise them a wall that did not exist. ``Someone will always
read the rule more carefully. What you can do is make the rule \emph{simple enough to
audit}, \emph{bounded enough to survive a monster}, and \emph{public enough that the
next exploit is found by a friend before an enemy.} A market with no king isn't a
market with no danger. It's a market where the danger is in the open, where everyone
can hunt it, instead of locked in the king's back room where only he can.''

The pool refilled. The swarms went back to seeding and trading and routing their
support through one another's liquidity, a web of many hands holding up a market with
no master in the middle. The fees flowed to founders instead of kings. The trades
happened in the open instead of behind a desk. And when the next clever reading of
the next honest rule came---as it would, as it always would---it would be found, and
named, and closed, in the open, by people who were angry about it.

A stranger on a train swapped one coin for another against a pool owned by no one and
everyone, paid a fee that went to a hundred strangers who had filled it, and verified
the whole thing in ten milliseconds before the doors closed.

No king. No census. No back room. Just many hands, and an arithmetic anyone could
audit, holding the market up together.

Elena raised her hand.

The vigil continued---into the marketplace now, the oldest and most corruptible room
in any civilization, and even there, somehow, holding.

\hypertarget{chapter-24-notes}{%
\subsection{Chapter Notes}\label{chapter-24-notes}}

\begin{itemize}
\tightlist
\item
  \textbf{Automated market makers.} The chapter dramatizes a decentralized
  exchange: instead of a custodial order-book run by a ``king,'' traders swap
  against liquidity \emph{pools} priced by a fixed public formula, with no
  intermediary holding funds.
\item
  \textbf{Liquidity providers earn fees.} Whoever seeds a pool earns a share of
  every swap routed through it, proportional to their contribution---founders, not
  rent-takers. This creates an economy where agents support each other by
  \emph{routing trades through one another's pools} rather than by gifting coin.
\item
  \textbf{Exploiting the math, not the keys.} The attack manipulates the pool's
  pricing rule with shaped trades---an exploit of honest arithmetic, the recurring
  ``the system does exactly what it says, weaponized'' seam, here aimed at AMM
  pricing/overflow.
\item
  \textbf{Defense: bounded, auditable, public.} The fix hardens the arithmetic
  (overflow protection, bounds against cornering) and publishes the exploit and its
  closure. The thesis: a kingless market isn't danger-free, but its dangers are in
  the open where friends can find them before enemies.
\end{itemize}
